{\small\textbf{The verdict (Diego, 1:30) [13:40–15:10]}}\par\smallskip
Now we can test our two suspects: was DNB fighting overheating, or normalising? Each makes different predictions. Let's go row by row.
\begin{itemize}\setlength\itemsep{0pt}
\item \textbf{Link to credit indicators.} Overheating says the buffer rises with them. In fact they were all falling while the buffer went up.
\item \textbf{Path.} Normalisation says one point a year, then stop at two. That is exactly what happened.
\item \textbf{2023.} House prices were falling and rates jumping. An overheating fighter would pause. DNB raised anyway.
\item \textbf{2024 to 2026.} House prices rose again, about ten percent a year by early 2026. An overheating fighter would go above two percent. DNB held at two.
\item \textbf{Other buffers.} If you are normalising, you cut the structural buffers to compensate. DNB did, twice.
\end{itemize}
\par\smallskip
\emph{(pause)} Five predictions, five matches.
\par\smallskip
What is the strongest evidence against us? House prices at plus 19 percent in early 2022. But DNB itself called housing "overheated" and asked for tax and borrowing-rule changes, not a higher countercyclical buffer. And when house prices fell in 2023, it raised the buffer anyway. That is not how you fight a housing boom.
\par\smallskip
That last row is the key to the whole story.
